\section{Julia and Rust alternatives for instruments and discontinuities}
\label{exec23:otheriv-scope}

A credit eligibility threshold produces several different empirical
questions. The discontinuity in loan approval can differ from the
discontinuity in loan receipt, and both differ from the effect of receiving
credit on a later monetary outcome. A bank comparing software must preserve
the assignment rule, the observation clock, the local population and the
uncertainty of every estimated discontinuity. Language coverage is useful
only when it exposes those distinctions. This section examines five concrete
Julia and Rust implementations at revisions preceding 7 October 2026. The
selected sources provide prospective alternatives and reasons for refusal;
they do not establish native numerical qualification or exhaustive coverage
of either language.
\citep{exec23_otheriv_source_0,exec23_otheriv_source_1,exec23_otheriv_source_2,exec23_otheriv_source_3,exec23_otheriv_source_4}

\subsection{Fixed effects and instrumental variables in Julia}
\label{exec23:otheriv-fe}
FixedEffectModels.jl, version 2.1.0 at the selected 30 September 2026
revision, provides a substantive fixed-effect and two-stage least-squares
route. Its formula distinguishes exogenous regressors, endogenous
regressors, instruments and absorbed effects. Let $X$ contain the structural
regressors and $Z$ the included instruments after the declared weighted
fixed-effect transformation. In the full-rank case, the familiar projection
representation is
\begin{equation}
 P_Z=Z(Z'Z)^{-1}Z',\qquad
 \widehat\beta=(X'P_ZX)^{-1}X'P_ZY.
 \label{exec23:otheriv-2sls}
\end{equation}
This numerical formula presumes sufficient rank. A causal interpretation
additionally needs a defended instrument exogeneity and exclusion argument,
and the intended homogeneous effect or heterogeneous local-effect
interpretation. An instrument-count check cannot establish any of those
scientific claims.
\citep{exec23_otheriv_source_0}

The inspected implementation retains finite-input checks, positive-weight
sample selection, singleton and column removals, a relative sweep tolerance
and first-stage residuals. It can recategorize a numerically collinear
endogenous regressor as exogenous, with an informational message. The
adapter must report that change and its reason. It cannot silently transfer
the user's original instrument interpretation to a different model.
Coefficients omitted for collinearity are reinserted as zero with missing
covariance. That zero is a representation of omission, not an estimated
zero effect. \citep{exec23_otheriv_source_0}

Two further distinctions matter operationally. First, the fixed-effect
solver's convergence flag is retained, but nonconvergence is logged and
fitting continues. Returning a fitted object therefore does not establish
that the residualization reached its requested accuracy. Second, the
first-stage rank-test call is inside an exception handler; an unsuccessful
test leaves missing statistics. The proposed worker should distinguish a
weak statistic, an unavailable statistic and a failed calculation. The
reported rank diagnostic is not a weak-instrument-valid confidence set.
Actual FixedEffects and Vcov dependency artifacts require their own
numerical and inferential review.
\citep{exec23_otheriv_source_0}

The source rejects particular HC2/HC3 combinations with absorbed effects
and exposes cluster and fixed-effect degree adjustments. It also floors
the residual inference degrees of freedom at one. A wrapper must retain
the original cluster/sample counts and establish adequate support rather
than treating that numerical floor as evidence of information. The fitted
object stores a sample mask and formula schema but not the original data,
so exact reconstruction needs external custody of input rows and their
ordering. The package's MIT notice and Julia 1.10 compatibility are useful
planning information, not clearance of every dependency or proof of a
particular installed CPU/GPU profile.
\citep{exec23_otheriv_source_0}

The older InstrumentalVariables.jl is a different candidate, with a
selected August 2016 revision and a REQUIRE file declaring Julia 0.3. Its
source implements Cholesky-based IV projection and covariance methods,
but several local contracts make it unsuitable for direct adoption in
the proposed engine. Three dimension guards construct a
\texttt{DimensionMismatch} object without explicitly throwing it; the
instrument-count expression compares a regressor count with the number
of instrument rows. Later operations may still throw, but these guards
are not reliable evidence that invalid dimensions are rejected at the
intended boundary.
\citep{exec23_otheriv_source_2}

Its coefficient-table expression applies a normal upper tail to the
squared statistic, despite a two-sided label. This should be treated as
a source-expression concern requiring an appropriate legacy runtime to
determine whether evaluation succeeds. Its no-argument weighted covariance
method also retrieves the stored residual vector and multiplies it in
place by square-root weights. If invoked repeatedly on that mutable
object, a residual of one with weight four becomes two and then four; its
squared contribution changes from four to sixteen. This exact mutation
witness is independent of the rest of the covariance calculation. It
shows why repeated accessor calls must be tested for purity, not that an
upstream Julia execution has been performed here. The README's default
HC0 wording also needs reconciliation with the inspected default scalar
residual-variance covariance formula.
\citep{exec23_otheriv_source_2,exec23_otheriv_reference}

\subsection{Sharp RD, curvature bounds and probability conventions}
\label{exec23:otheriv-sharp}
RegressionDiscontinuity.jl declares version 0.4.0 at its selected May 2024
revision and explicitly warns that the package is experimental. Its
running-variable object records whether treatment is assigned above or
below the threshold, including equality. This is a valuable feature: a
credit policy applying at score $s\geq c$ is not the same data rule as
$s>c$ when many accounts lie exactly at $c$. The local-linear interface
fits a weighted outcome regression with treatment, centered score and
their interaction, and identifies the treatment coefficient as the
sharp-RD contrast. It labels its inference as naive and its bias as
unaccounted. That limitation should remain visible in the product.
\citep{exec23_otheriv_source_1}

The inspected local-linear table computes $1-\Phi(|z|)$ for its p-value
but uses the usual two-sided $\alpha/2$ critical value for its interval.
The density-test implementation uses $2[1-\Phi(|z|)]$. At $z=2$, the
one-tail quantity is approximately $0.02275013$, the conventional
two-sided quantity is $0.04550026$, and the legacy IV expression
$1-\Phi(z^2)$ is approximately $0.00003167$. These are plainly different
probability conventions. A common output label does not reconcile them;
the adapter needs a declared hypothesis and a validated calculation.
The values are independently computed normal tails, not observed native
outputs. \citep{exec23_otheriv_source_1,exec23_otheriv_source_2,exec23_otheriv_reference}

The package also exposes an optimized sharp-RD route with an explicit
bound $B$ on the second derivative, a chosen discretization and an
external optimization solver. That is a distinct scientific contract
from bias-unadjusted local linear fitting. Its source constructs a
cone problem, calls optimization and then extracts solution values
without a local termination-status guard. A worker must check solver
status, primal/dual validity and numerical accuracy before treating the
reported maximum bias as a certificate. The curvature bound needs a
scientific justification and sensitivity analysis. The named paper and
the complete solver/dependency argument require their own original-source
review before any theorem-level guarantee is adopted.
\citep{exec23_otheriv_source_1}

Noise-induced randomization is another exposed research route, with
outcome bounds, a latent convex class, target, noise model, solver,
discretization and a switch that skips bias calculation. The presence
of that interface does not solve arbitrary noisy credit scores. A
proposal should preserve these assumptions and identify the extra
measurement information needed to justify them. Only selected entry
and optimization spans were inspected here; the full method, proofs
and empirical performance remain separate obligations.
\citep{exec23_otheriv_source_1}

\subsection{A separate Julia port of the rdrobust workflow}
\label{exec23:otheriv-port}
The repository xiangao/RDRobust.jl declares version 0.1.1 at its
3 October 2026 revision. Its repository identity and license notice
must be retained separately from any metadata author string or claimed
similarity to the R package. No official affiliation or endorsement is
inferred here. The README claims a numerical tolerance on selected test
cases, but those tests have not been executed in this review. The
selected source does expose substantive local-polynomial machinery: main
and bias bandwidths, covariates, clustering, nearest-neighbor residual
estimation, mass-point adjustments and sharp/fuzzy targets.
\citep{exec23_otheriv_source_3}

Its output distinguishes conventional, bias-corrected and robust rows.
The bias-corrected row uses the corrected point estimate with the
conventional standard error, while the robust row uses the corrected
estimate and its robust standard error. A common reporting layer must
not combine one row's coefficient with another row's interval without
an explicit reason. The fuzzy path carries a two-component ratio
gradient into its covariance calculation. That provides a useful
source-level contrast with the Rust implementation considered below,
while leaving the actual statistical and numerical qualification open.
\citep{exec23_otheriv_source_3}

Several implementation details need deliberate controls. The helper
named \texttt{qrXXinv} directly inverts a cross-product matrix. Its
complete-case helper excludes NaN but does not itself exclude infinity.
Cluster labels are converted to floating-point values, so an adapter
needs a lossless category map instead of blindly casting identifiers.
The cluster correction contains $G/(G-1)$ and $(n-1)/(n-k)$, requiring
valid cluster and residual degrees. When covariates become collinear
with local polynomials, the code warns and still uses a generalized
inverse. A computed number must not become an accepted identified
contrast merely because that inverse exists.
\citep{exec23_otheriv_source_3}

Mass-point counting and bandwidth enlargement help describe the local
sample; they do not create variation that is absent. Five observations
at score $-1$ on one side and five at score $1$ on the other still give
only one distinct score on each side. The two-column intercept/score
design has determinant zero within each side. Adding more duplicate
accounts does not identify a local slope. The response should report
support, rank and the assumed smoothness or discrete-score interpretation
explicitly. \citep{exec23_otheriv_source_3,exec23_otheriv_reference}

\subsection{Fuzzy ratios and a Rust example that should not be adopted}
\label{exec23:otheriv-ratio}
The selected Rust repository suenot/127-regression-discontinuity
contains a trading-oriented library named rdd_trading, version 0.1.0
at its March 2026 revision. It implements separate local fits and a
fuzzy ratio, but the inspected uncertainty and diagnostic behavior
do not support adoption as the proposed scientific engine. No trading,
exchange, network or downloaded-code execution is involved in this
review. A financial example in a repository is not evidence that
threshold crossing is a valid causal assignment mechanism.
\citep{exec23_otheriv_source_4}

Let $a$ denote an estimated outcome discontinuity and $b$ an estimated
treatment-receipt discontinuity. The fuzzy target is $\theta=a/b$.
Under a justified regular joint limit, with $b$ bounded away from zero,
the first-order variance uses
\begin{equation}
 g=\left(\frac1b,-\frac{a}{b^2}\right)',
 \qquad V_\theta\simeq g'V_{a,b}g.
 \label{exec23:otheriv-delta}
\end{equation}
This expression is not weak-instrument-robust inference. It explains
why dividing an outcome standard error by $|b|$ is generally incomplete:
that shortcut freezes the estimated denominator and omits its covariance
with the numerator. The inspected Rust fuzzy method uses that shortcut
and rejects first-stage magnitudes below a fixed $0.01$. A fixed numerical
cutoff is not a calibrated test of weak identification.
\citep{exec23_otheriv_source_4}

For a reproducible stipulated example, take
\begin{equation}
 a=2,\qquad b=\frac15,\qquad
 V_{a,b}=\begin{pmatrix}1/25&1/250\\1/250&1/100\end{pmatrix}.
 \label{exec23:otheriv-example}
\end{equation}
The ratio is ten dollars and $g=(5,-50)'$. The full delta variance is
$1+25-2=24$ dollars squared; freezing the denominator gives only one.
With an assumed six-dollar cost per additional treated account, the
ratio-based net is four dollars. This is an illustrative regular-case
calculation; the numbers do not establish that a small first stage
satisfies the regular approximation, that the instrument is valid, or
that a real bank earns this return. If the denominator is weak, the
appropriate route is a justified joint confidence-set or test-inversion
procedure, retaining unbounded or disconnected sets where warranted,
rather than adjusting a symmetric interval to look precise.
\citep{exec23_otheriv_reference}

The Rust sharp and fuzzy methods also set the test statistic to zero
when the computed standard error is nonpositive. Its weighted local
solver adds a fixed determinant ridge and floors variance denominators.
These choices need numerical and statistical justification; an interval
or p-value produced after such a fallback is not automatically usable.
The method named for an optimal bandwidth uses a simplified curvature
calculation, regularization and range bounds. Its label alone does not
establish correspondence to a named published selector.
\citep{exec23_otheriv_source_4}

Its density diagnostic is particularly important for banking use. The
source calculates near-cutoff counts and their density ratio. If either
local density is nonpositive, it substitutes a zero log contrast and a
standard error of one, yielding a p-value of one. The correct product
interpretation is unavailable local evidence, not reassurance that
scores were not manipulated. The binned counts calculated earlier in
the function are not used in that returned statistic. A declared
density-test name cannot substitute for inspection of the actual
calculation, and a failure to reject manipulation never certifies all
RD identification assumptions.
\citep{exec23_otheriv_source_4,exec23_otheriv_reference}

\subsection{Choosing a route and preserving omissions}
\label{exec23:otheriv-choice}
The useful proposal choice is method-specific. FixedEffectModels.jl is
a conditional Julia worker candidate after numerical, covariance and
failure-contract qualification. RegressionDiscontinuity.jl remains a
research reference whose naive and optimized/noise-based paths need
separate treatment. The old InstrumentalVariables.jl remains historical
reference material; its runtime and local hazards argue against current
native adoption. The distinct RDRobust.jl port is a comparator requiring
provenance, rank and matched-output trials. The inspected Rust trading
implementation is rejected as the current scientific engine because
its ratio uncertainty and unavailable-density fallbacks would require
substantive repair. All five retained license texts declare MIT terms,
but the actual numerical, optimization and runtime dependencies still
need distribution-specific review.
\citep{exec23_otheriv_source_0,exec23_otheriv_source_1,exec23_otheriv_source_2,exec23_otheriv_source_3,exec23_otheriv_source_4}

An integration trial should compare the complete intervention request
under declared assumptions, not just regression coefficients. Retain
sample exclusions, score support, timing, first-stage strength, complete
covariance, numerical convergence, diagnostics and failures. Use the
independent ratio, probability, rank and mutation witnesses as initial
oracles, followed by predeclared public or synthetic calibration and
cost trials in an authorized installed environment. The scientific
owner approves the estimand and inference; the maintainer owns exact
versions and regression cases; the operator owns custody and complete
cost accounting. A new port needs a demonstrated benefit before
duplicating these obligations. \citep{exec23_otheriv_reference}

Eight additional Julia repositories returned by the bounded discovery
queries remain uninspected, including alternative IV, shift-share,
system-projection and model-averaging routes. Other languages and
registries are also unresolved. A zero-result Rust IV query is a
coverage limit, not proof that no implementation exists. These omissions
prevent a claim of comprehensive coverage while allowing the concrete
source findings and proposed comparisons above to guide the next
prototype decision.
